\clearpage\begingroup\SourceLeaf{348}\ArithmeticHead{17}{R}{DIE GRUNDLAGEN DER GEOMETRIE.}
\begin{center}
\textls[100]{A. Von den Anfängen der nichteuklidischen Geometrie bis zur}\\
\textls[100]{Entdeckung der transzendenten Trigonometrie (1792—1817).}\par\medskip
\textbf{3.}\par\smallskip
\textls[100]{Einleitendes. Die Jugendzeit (1792—1795).}
\end{center}
\noindent\hspace*{\parindent}Als \textsc{Jacobi} am 5. August 1827, auf Grund eines Briefes, den \textsc{Schumacher}
an ihn gerichtet hatte, \textsc{Legendre} mitteilte, \textsc{Gauss} habe schon 1808 einen Teil
der von \textsc{Jacobi} in den Astronomischen Nachrichten veröffentlichten Sätze besessen
(\textsc{Jacobi}, Werke I, S. 394), antwortete \textsc{Legendre} am 30. November:
»Comment se fait-il que M. \textsc{Gauss} ait osé vous faire dire que la plupart de
vos théorèmes lui étaient connus et qu’il en avait fait la découverte dès 1808?
Cet excès d’impudence n’est pas croyable de la part d’un homme qui a assez
de mérite personnel pour n’avoir pas besoin de s’approprier les découvertes des
autres« (S. 398), und am 14. April 1828 setzt er hinzu: »Il y a des gens comme
M. \textsc{Gauss}, qui ne se feraient pas scrupule de vous ravir, s’ils le pouvaient,
le fruit de vos recherches, et de prétendre qu’elles sont depuis longtemps en
leur possession. Prétention bien absurde assurément: car si M. \textsc{Gauss} était
tombé sur de pareilles découvertes qui surpassent, à mes yeux, tout ce qui a
été fait jusqu’ici en analyse, bien sûrement il se serait empressé de les publier«
(S. 418).
\par
Der Nachlass von \textsc{Gauss} hat demgegenüber gezeigt, dass dieser bereits
im Jahre 1797 begonnen hatte, die lemniskatischen Funktionen zu untersuchen,
dass er bis zum Jahre 1800 die wesentlichen Eigenschaften der allgemeinen
elliptischen Funktionen erkannt hatte und dass er im Jahre 1808 diese Untersuchungen
wieder aufgenommen und sich dem Problem der Teilung zugewandt
hatte, auf das sich jene von \textsc{Jacobi} entdeckten Sätze beziehen.
\par
Ebenso sind in anderen Fällen die Angaben, die \textsc{Gauss} über seine mathematischen
Entdeckungen gemacht hat, durch Aufzeichnungen im Nachlass oder
durch Briefe bis in die Einzelheiten hinein bestätigt worden. Wie konnte es
auch anders sein bei einem Manne von so grosser Wahrheitsliebe und Gewissenhaftigkeit?
Dazu wurde \textsc{Gauss} durch ein ungewöhnlich treues Gedächtnis
unterstützt. Auch hat er häufig die Aufzeichnungen aus den Jahren 1796
bis 1815 benutzt, die er sich in einem Notizenjournal oder \textit{Tagebuch} gemacht
hatte (W. X\textsubscript{1}, S 488—572); in der späteren Zeit pflegte er umgekehrt in Handbücher
\par\vfill\noindent{\footnotesize X\textsubscript{2} Abh. 4.\hfill 3\hspace*{2em}}
\clearpage\SourceLeaf{349}\ArithmeticHead{18}{L}{STÄCKEL, GAUSS ALS GEOMETER.}
\noindent kurze Bemerkungen über mathematische Sätze einzutragen, die er in
Briefen erwähnt hatte. Gewiss kommen gelegentlich Angaben vor, die einander
zu widersprechen scheinen, allein in den allermeisten Fällen haben sie
sich bei sorgfältiger Deutung in Übereinstimmung bringen lassen, und so wird
man den Äusserungen von \textsc{Gauss} über die Entstehung seiner Gedanken volles
Vertrauen entgegenbringen dürfen.
\par
Hiernach sind auch die Äusserungen zu beurteilen, die \textsc{Gauss} über die
Anfänge seiner Beschäftigung mit den Grundlagen der Geometrie gemacht hat.
\par
Am 28. November 1846 schreibt \textsc{Gauss} an \textsc{Schumacher}, er habe schon
im Jahre 1792, also mit 15 Jahren, an eine Geometrie gedacht, »die stattfinden
müsste und strenge konsequent stattfinden könnte, wenn die \textsc{Euklid}ische
Geometrie nicht die wahre ist«, das heisst, wenn das elfte Axiom nicht gilt
(W. VIII, S. 238). Hiermit ist jedenfalls nur das erste Aufblitzen des Gedankens
gemeint. Denn unmittelbar vorher, am 2. Oktober 1846, hatte \textsc{Gauss}
zu \textsc{Gerling} geäussert, der Satz, dass in jeder vom Parallelenaxiom unabhängigen
Geometrie der Flächeninhalt eines Vielecks der Abweichung der
Summe der Aussenwinkel von \(360^0\) proportional ist, sei »der erste, gleichsam
an der Schwelle liegende Satz der Theorie, den ich schon im Jahr 1794 als
notwendig erkannte« (W. VIII, S. 266). Wir werden sehen, dass diese Beziehung
zwischen dem Inhalt und der Winkelsumme eines Vielecks einen
Angelpunkt der \textsc{Gauss}schen Theorie gebildet hat, und dürfen daher annehmen,
dass der Zeitpunkt, wo er zu einer solchen grundlegenden Einsicht gelangt
war, sich ihm fest eingeprägt hatte.
\par
Als \textsc{Wolfgang Bolyai} seinem Jugendfreunde die \textit{Scientia spatii absolute
vera} seines Sohnes \textsc{Johann} übersandt hatte, bemerkte \textsc{Gauss} am 6. März
1832, der ganze Inhalt der Schrift komme fast durchgehends überein »mit
seinen eigenen, zum Teile schon seit 30 bis 35 Jahren angestellten Meditationen«
(W. VIII, S. 221). Man wird damit bis auf die Jahre von 1797 bis
1802 zurückgeführt. Zu dieser Zeit hat \textsc{Gauss} also angefangen, in weiterem
Umfange die Folgen zu entwickeln, die sich ergeben, wenn man die Wahrheit
des elften \textsc{Euklid}ischen Axioms leugnet. In der Tat bringt das \textit{Tagebuch}
unter dem September 1799 (T. Nr. 99) die Eintragung: »In principiis geometriae
egregios progressus fecimus«. Worin diese ausgezeichneten Fortschritte
bestanden haben, wird noch zu erörtern sein.
\clearpage\SourceLeaf{350}\ArithmeticHead{19}{R}{DIE GRUNDLAGEN DER GEOMETRIE.}
\noindent\hspace*{\parindent}Gehen wir in der Reihe der Zeugnisse weiter. Kurz vorher, am 17. Mai
1831, hatte \textsc{Gauss} an \textsc{Schumacher} berichtet, er habe »angefangen, einiges von
seinen Meditationen über die Parallellinien aufzuschreiben, die zum Teil schon
gegen 40 Jahr alt sind« (W. VIII, S. 213). Er geht also hier bis auf die
keimhaften Ursprünge zurück, für die er die Jahre 1792 und 1794 genannt
hatte. Dieselbe Datierung findet sich in dem Briefe an \textsc{Taurinus} vom 8. November
1824: »Ich vermute, dass Sie sich noch nicht lange mit diesem Gegenstande
[der Parallelentheorie] beschäftigt haben. Bei mir ist es über 30 Jahr,
und ich glaube nicht, dass jemand sich eben mit diesem zweiten Teil [wo die
Winkelsumme des Dreiecks kleiner als zwei Rechte ist] mehr beschäftigt haben
könne als ich, obgleich ich niemals darüber etwas bekannt gemacht habe«
(W. VIII, S. 186).
\clearpage\endgroup
